(1) 约化坐标在接触阶段满足 \(\mu\ddot\delta=-K\delta^{3/2}\)，初始 \(\delta(0)=0,\dot\delta(0)=v\)。乘 \(\dot\delta\) 并积分得 \(\boxed{\frac\mu2\dot\delta^2+\frac{2K}{5}\delta^{5/2}=\frac\mu2v^2}\)。非接触阶段压入量被限制为0，不能让Hertz力扩展至负压入量。

(2) 最大压缩处 \(\dot\delta=0\)，故 \(\boxed{\delta_{\max}=(5\mu v^2/(4K))^{2/5}}\)、\(\boxed{F_{\max}=K\delta_{\max}^{3/2}}\)。因此 \(\delta_{\max}\propto v^{4/5}K^{-2/5}\)、\(F_{\max}\propto v^{6/5}K^{2/5}\)，与线性弹簧接触的指数不同。

(3) 由第一问 \(\dot\delta=v\sqrt{1-(\delta/\delta_{\max})^{5/2}}\)，压缩时间 \(t_-=\delta_{\max}v^{-1}\int_0^1[1-u^{5/2}]^{-1/2}du\)。取 \(z=u^{5/2}\) 得 \(t_-=(2\delta_{\max}/(5v))B(2/5,1/2)\)。回弹过程时间反演，故 \(\boxed{t_c=(4\delta_{\max}/(5v))B(2/5,1/2)}\)，并有 \(t_c\propto\mu^{2/5}K^{-2/5}v^{-1/5}\)。

(4) 接触力在压缩和回弹阶段均为排斥力。分离时压入量归零，守恒第一积分给相对速度 \(-v\)，因此恢复系数 \(e=1\)；总相对动量改变量的冲量大小 \(\boxed{J=2\mu v}\)。对初速度 \(v_1,v_2\)、法向相向且 \(v_1-v_2=v\)，末速度满足 \(v_1'=v_1-J/m_1\)、\(v_2'=v_2+J/m_2\)，总动能无损失；实际黏弹材料需另加耗散力。
